% TOP_PROBLEM: {"id": 2587, "title": "Kourovka Notebook Problem 21.78", "category_id": 17, "category": {"id": 17, "name": "group_theory", "display_name": "Group Theory", "description": "Problems about groups, group actions, representations, and related algebraic structures.", "slug": "group-theory", "order_index": 17, "created_at": "2026-07-31T15:26:25.671Z"}, "status": "open", "problem_number": "KOU-21.78", "set_id": 10}
% FIRST_POSTED: 2026-10-01
% ENTRY_KIND: statement_only
% UnsolvedMath ID 2587; source catalogue code KOU-21.78
% !TeX program = lualatex
% Definition and sources only; this file is not an LLM solution attempt.
\documentclass[11pt,a4paper]{article}
\usepackage[margin=25mm]{geometry}
\usepackage{fontspec}
\setmainfont{lmroman10-regular.otf}[BoldFont=lmroman10-bold.otf,
 ItalicFont=lmroman10-italic.otf,BoldItalicFont=lmroman10-bolditalic.otf]
\usepackage{amsmath,amssymb,mathtools}
\usepackage{unicode-math}
\setmathfont{latinmodern-math.otf}
\AtBeginDocument{\let\setminus\smallsetminus}
\usepackage{xurl,hyperref}
\hypersetup{colorlinks=true,urlcolor=blue}
\setcounter{secnumdepth}{0}
\setlength{\parindent}{0pt}
\setlength{\parskip}{5pt}
\setlength{\emergencystretch}{3em}
\begin{document}
\section*{Kourovka Notebook Problem 21.78}\label{2587}
{\small UnsolvedMath ID 2587 · KOU-21.78 · Group Theory · Kourovka Notebook - New Problems, Issue 21}
\subsection{Definitions and mathematical statement}
Fix a prime $p$. A pro-$p$ group is a topological group isomorphic to an inverse limit of finite groups of $p$-power order, with their discrete topologies. Give $\mathbb F_p$ the discrete topology and the trivial $G$-action. Its continuous cochain groups are
\[
 C^n(G,\mathbb F_p)=\{\hbox{continuous maps }G^n\longrightarrow\mathbb F_p\}.
\]
Here $C^0=\mathbb F_p$. With trivial coefficients the differential is
\[
\begin{split}
(df)(g_1,\ldots,g_{n+1})={}&f(g_2,\ldots,g_{n+1})\\
&+\sum_{i=1}^n(-1)^if(g_1,\ldots,g_ig_{i+1},\ldots,g_{n+1})\\
&+(-1)^{n+1}f(g_1,\ldots,g_n).
\end{split}
\]
The continuous cohomology $H^n(G,\mathbb F_p)$ is $\ker d/\operatorname{im}d$ in degree $n$. Cup product, induced on cochains of degrees $r,s$ by
\[
 (f\smile h)(g_1,\ldots,g_{r+s})
 =f(g_1,\ldots,g_r)h(g_{r+1},\ldots,g_{r+s}),
\]
makes $H^*(G,\mathbb F_p)=\bigoplus_{n\ge0}H^n(G,\mathbb F_p)$ a graded algebra.

Suppose that $H^n(G,\mathbb F_p)$ is finite for every $n\ge0$, equivalently finite-dimensional over $\mathbb F_p$ in each degree. Must this graded algebra be finitely generated over $\mathbb F_p$? Precisely, must there be finitely many homogeneous classes $u_1,\ldots,u_t$ such that every cohomology class is an $\mathbb F_p$-linear combination of products of these classes and the unit? The finite generating list may depend on $G$ and $p$; the hypothesis gives no common dimension bound across degrees.
\subsection{Short English statement}
Does finiteness of every continuous cohomology group of a pro-$p$ group force finite generation of its whole mod-$p$ cohomology algebra?
\subsection{Sources}
\begin{itemize}
\item[S1] ULAM.AI and the UnsolvedMath Contributors, supplied problems.json, record 2587 (KOU-21.78), Kourovka Notebook - New Problems, Issue 21. Catalogue identity and source transcription; CC BY 4.0.
\url{https://huggingface.co/datasets/ulamai/UnsolvedMath}.
\item[S2] The Kourovka Notebook, 21st edition, version 47 (30 September 2026), new problems, Problem 21.78. Expanded formulation of the supplied catalogue statement.
\url{https://arxiv.org/pdf/1401.0300v47}.
\end{itemize}
\end{document}
